\section{Ideas básicas del aprendizaje por imitación}

El aprendizaje por imitación (Imitation Learning) parte de copiar directamente al experto, en lugar de diseñar una función de recompensa compleja. Es adecuado para escenarios con datos suficientes y demostraciones de alta calidad del experto, pero donde el reward no está bien definido.

\begin{importantbox}{El objetivo del aprendizaje por imitación}
Dado un conjunto de trayectorias de demostración $\mathcal{D} = \{(s_1,a_1,\ldots,s_T)\}$ generado por una política experta $\pi_{\text{expert}}$, el objetivo es encontrar una política parametrizada $\pi_\theta$ cuya distribución de acciones predicha coincida lo más posible con la del experto en los mismos estados, de modo que se conserve un desempeño al nivel del experto.
\end{importantbox}

\subsection{Version 0: Regresión de política determinista}

El enfoque más directo es tratar el aprendizaje por imitación como un mapeo de comportamiento:
$$\min_\theta \frac{1}{|\mathcal{D}|}\sum_{(s,a) \in \mathcal{D}} \|\pi_\theta(s) - a\|^2.$$
Este método depende por completo de que los datos de demostración estén bien balanceados y es muy sensible al noise.

\begin{warningbox}{Problema de la media y conductas irreversibles}
Cuando los datos del experto presentan multimodalidad (por ejemplo, girar a la izquierda o a la derecha), el modelo de regresión produce una acción «promedio», lo cual provoca que el sistema entre en estados peligrosos dentro del entorno; este problema es habitual en tareas como la conducción urbana y la manipulación con brazos dobles.
\end{warningbox}

\subsection{Clonación de comportamiento vs. aprendizaje por refuerzo inverso}

La clonación de comportamiento modela directamente las acciones, mientras que el aprendizaje por refuerzo inverso (Inverse RL) reconstruye el reward y después realiza la optimización mediante RL. A continuación se compara el trade-off entre ambos enfoques:

\begin{table}[H]
\centering
\begin{tabular}{p{0.25\textwidth} p{0.32\textwidth} p{0.32\textwidth}}
\toprule
Método & Ventajas & Desventajas \\
\midrule
Clonación de comportamiento & Entrenamiento sencillo, usa directamente supervised learning & Pierde fácilmente la multimodalidad y se ve muy afectada por el cambio de distribución \\
RL inverso & Puede recuperar el reward latente y cuenta con capacidad de generalización & Requiere una gran cantidad de muestreo; el aprendizaje del reward está acoplado con el entrenamiento de la policy \\
\bottomrule
\end{tabular}
\caption{Comparación entre la clonación de comportamiento y el RL inverso}
\end{table}

\begin{knowledgebox}{Por qué priorizar la clonación de comportamiento}
En muchos escenarios industriales (como la conducción autónoma), las trayectorias del experto son suficientes y de alta calidad; en ese caso, aplicar directamente la clonación de comportamiento reduce en gran medida el ciclo de lanzamiento. Si aparece deriva de la distribución, se introduce gradualmente un modelo de reward.
\end{knowledgebox}

\subsection{Puesta en producción en escenarios reales y caution}

Las estrategias habituales en el despliegue real incluyen: el sobremuestreo de los datos de conducción, entrenar una fallback policy usando a los experts como baseline, y establecer un mecanismo de monitoreo human-in-the-loop. En escenarios de fallas de baja frecuencia, es un diseño habitual usar primero la clonación de comportamiento para hacer un warm-start y después aplicar fine-tuning con RL.

\begin{importantbox}{La interpretabilidad de la clonación de comportamiento}
Debido a que la salida es deterministic/log-prob, resulta fácil hacer trace-back; en cuanto el modelo produce una salida anómala, el equipo de ingeniería puede revisar directamente el expert sample más reciente para ubicar rápidamente el origen del problema.
\end{importantbox}

%-\subsection{Resumen de la sección}
-
%-La clonación de comportamiento ofrece...
%-\subsection{Fusión de múltiples expertos y pesos dinámicos}
-
-En la era de los grandes modelos...

\begin{knowledgebox}{El valor de la mezcla de expertos}
Colocar los logits de múltiples expertos dentro de un marco Mixture-of-Experts permite lograr una transferencia zero-shot con pocas muestras, porque la gating network puede orientarse automáticamente hacia el expert más relevante para la tarea actual.
\end{knowledgebox}

\subsection{Resumen de la sección}

La clonación de comportamiento ofrece la ruta de entrada más rápida, pero la regresión deterministic debe combinarse con gating de múltiples expertos y modelado de distribución para poder generalizar en escenarios complejos. La mezcla de múltiples expertos ofrece una asignación automática entre distintas skill, lo que amplía los límites de aplicabilidad del aprendizaje por imitación.

